\section{Введение}
\subsection*{Цель работы} \addcontentsline{toc}{subsection}{Цель работы}

Изучить алгоритмы оценки среднего значения эргодического случайного процесса, исследовать влияние корреляции между отсчётами на погрешность оценки и освоить способы оценки среднего, дисперсии и доверительного интервала по временным и спектральным характеристикам сигнала.

\subsection*{Объект и задачи исследования} \addcontentsline{toc}{subsection}{Объект и задачи исследования}

Объект исследования --- дискретный гауссовский шум, генерируемый лабораторной установкой при относительных временах корреляции $\tau_\text{корр} = 10$, $30$ и $100$. В работе решаются следующие задачи:

\begin{itemize}
	\item наблюдение реализаций шума и СПМ при разных $\tau_\text{корр}$;
	\item исследование оценки среднего, её разброса и зависимости СКО от $N$ и $\Delta t$;
	\item определение характеристик исходного и усреднённого процессов по СПМ;
	\item исследование гистограмм оценок и доверительных интервалов при разных $N$ и $\beta$.
\end{itemize}
